% ------------------------------------------------------------
% A6.x Substrate Uplift Module E — Economic Trajectory Invariants (v0.1)
% Author: Borealis S. Hedling
% Altitude: A6.x (Substrate Uplift Layer)
% Purpose: Define the trajectory-invariant substrate required
%          for semantic stability, provenance stability, modal
%          continuity, and collapse-compatible trajectory geometry.
% ------------------------------------------------------------

\documentclass[11pt]{article}
\usepackage{amsmath, amssymb, amsthm}

\title{A6.x Substrate Uplift Module E: Economic Trajectory Invariants}
\author{Borealis S. Hedling}
\date{October 2026}

\begin{document}
\maketitle

\section*{Overview}

This module defines the economic trajectory invariant substrate required for
A6 altitude geometry. Trajectory invariants provide the stability and continuity
structure required by the A6 Constructive Fixes and the A7 Composite Substrate
Operator Algebra.

% ------------------------------------------------------------
\section{Trajectory Definition}

Let $\gamma : [0,1] \to M$ be a smooth trajectory.

\textbf{Definition (Trajectory).}
\[
\gamma(t) \in M, \quad \gamma(0) = p, \quad \gamma(1) = q.
\]

% ------------------------------------------------------------
\section{Economic Trajectory Invariant}

Define the economic invariant:
\[
I_{\mathrm{econ}}(\gamma) = \int_0^1 \Phi(\gamma(t)) \, dt,
\]
where $\Phi$ is a potential function.

\textbf{Theorem (Invariance).}
If $\Phi$ is smooth, then:
\[
I_{\mathrm{econ}}(\gamma) \text{ invariant under reparameterization}.
\]

This enforces Fix 12 (semantic lift preservation).

% ------------------------------------------------------------
\section{Trajectory Stability}

Define trajectory stability:
\[
S(\gamma) = \sup_{t \in [0,1]} \|\dot{\gamma}(t)\|.
\]

\textbf{Theorem (Stability Bound).}
If $S(\gamma) < \sigma_{\max}$, then:
\[
\gamma \text{ stable}.
\]

This enforces Fix 9 (stability envelope monotonicity).

% ------------------------------------------------------------
\section{Modal Trajectory Lift}

Let $\mu$ be the modal map.

\textbf{Definition (Modal Lift).}
\[
\mu(\gamma(t)) \leadsto \mu(\gamma(t+\epsilon)).
\]

\textbf{Theorem (Lift Preservation).}
If $\gamma$ is stable, modal truth lifts along $\gamma$.

This enforces Fix 10 (modal lift preservation).

% ------------------------------------------------------------
\section{Provenance Trajectory Lift}

Define provenance lift:
\[
I_{\mathrm{prov}}(\gamma(t)).
\]

\textbf{Theorem (Provenance Stability).}
If $\gamma$ is stable, provenance lifts along $\gamma$.

This enforces Fix 11 (provenance lift preservation).

% ------------------------------------------------------------
\section{Collapse-Surface Compatibility}

Let $\Sigma$ be the collapse surface.

\textbf{Definition (Collapse-Compatible Trajectory).}
\[
\gamma(t) \in \Sigma \Rightarrow S(\gamma) = 0.
\]

\textbf{Theorem (Local Detectability).}
If $\gamma$ intersects $\Sigma$, collapse is locally detectable.

This enforces Fix 8 (collapse locality).

% ------------------------------------------------------------
\section{Spectral Trajectory Compatibility}

Define:
\[
\mathcal{S}_{\gamma}(t) = \mathrm{Spec}(L(\gamma(t))).
\]

\textbf{Theorem (Spectral Continuity).}
If $\gamma$ is stable, then:
\[
\mathcal{S}_{\gamma}(t) \text{ continuous in } t.
\]

This enforces Fix 5 (spectral continuity).

% ------------------------------------------------------------
\section{Conclusion}

The economic trajectory invariant substrate defined here provides the semantic,
provenance, modal, collapse-compatible, and spectral stability structures
required for A6 Constructive Fixes and A7 Composite Substrate Operator Algebra.
It is the fifth foundational substrate module in the A6.x uplift sequence.

\end{document}
